\documentclass[11pt]{article}

\usepackage[margin=1in]{geometry}
\usepackage{amsmath,amssymb,amsthm}
\usepackage{hyperref}
\usepackage{microtype}

\title{\textbf{Admissible Mixed Hodge Fixed-Part}}
\author{Inacio F.\ Vasquez \\ Independent Researcher}
\date{}

% --- Theorem Environments ---
\theoremstyle{plain}
\newtheorem{theorem}{Theorem}
\newtheorem{lemma}[theorem]{Lemma}
\newtheorem{corollary}[theorem]{Corollary}

\theoremstyle{remark}
\newtheorem{remark}[theorem]{Remark}

% --- Notation ---
\newcommand{\Q}{\mathbb{Q}}
\newcommand{\mathbbV}{\mathbb{V}}

\begin{document}

\maketitle

\begin{center}
\textbf{Status summary.}
\begin{itemize}
\item Fixed-part and horizontal subtori results proven for admissible VMHS.
\item Dependencies: admissibility, functorial filtrations, Deligne canonical extension.
\item Lean references: \texttt{URF\_Core.lean → Hodge.*}
\end{itemize}
\end{center}

\tableofcontents

\section{Monodromy Invariants}

\begin{lemma}[Monodromy invariants]\label{lem:hodge-monodromy}
For an admissible $\mathbbV$, flat sections invariant under the Gauss--Manin connection form a sub-VMHS over $\Q$.
\end{lemma}

\begin{proof}
Deligne canonical extension provides compatible weight/Hodge filtrations.  
Flat sections respect these filtrations and are invariant under connection.  
\textbf{Lean stub reference:} \texttt{URF\_Core.lean → Hodge.monodromy\_invariants}
\end{proof}

\section{Fixed Part Theorem}

\begin{theorem}[Fixed part for admissible VMHS]\label{thm:hodge-fixed-part}
Flat sections $H^0(S, \mathbbV_\Q)$ form a rational mixed Hodge substructure.
\end{theorem}

\begin{proof}
Admissibility ensures well-behaved filtrations.  
Functoriality under Gauss--Manin connection ensures rationality.  
\textbf{Lean stub reference:} \texttt{URF\_Core.lean → Hodge.fixed\_part}
\end{proof}

\section{Rationality of Horizontal Subtori}

\begin{corollary}[Rational horizontal subtori]\label{cor:hodge-rational}
If $T$ is a horizontal real subtorus generated by an admissible normal function, then $T$ is rational: $T = \mathsf{Flip}(T)$.
\end{corollary}

\begin{proof}
Apply Proposition (Flip preserves horizontality) and Theorem \ref{thm:hodge-fixed-part}.  
The minimal rational envelope forces $T = \mathsf{Flip}(T)$.  
\textbf{Lean stub reference:} \texttt{URF\_Core.lean → Hodge.rational\_horizontal\_subtori}
\end{proof}

\begin{remark}[Scope]
All results restricted to admissible VMHS.  
\textbf{Frozen module:} Yes.
\end{remark}

\end{document}

